\documentclass[12pt,a4paper]{article}
\usepackage[UTF8]{ctex}
\usepackage{amsmath,amssymb,amsthm}
\usepackage{geometry}

\geometry{margin=2.5cm}

\title{高度齿轮传动机器人的导纳控制与刚性环境模拟}
\author{第11章补充说明}
\date{}

\begin{document}

\maketitle

\section{问题提出}

\textbf{问题}：为什么高度齿轮传动的机器人的导纳控制可以很好地模拟刚性环境？

\textbf{原文}：
\begin{quote}
模拟低质量环境对于导纳控制的机器人来说是具有挑战性的，因为小的力产生大的加速度。有效的大增益可能产生不稳定。然而，高度齿轮传动的机器人的导纳控制可以很好地模拟刚性环境。
\end{quote}

\section{高度齿轮传动的特点}

\subsection{齿轮传动的原理}

\textbf{齿轮比（Gear Ratio）}：
\begin{equation}
N = \frac{\omega_m}{\omega_j} = \frac{\tau_j}{\tau_m}
\end{equation}

其中：
\begin{itemize}
\item $\omega_m$：电机角速度（输入）
\item $\omega_j$：关节角速度（输出）
\item $\tau_m$：电机力矩（输入）
\item $\tau_j$：关节力矩（输出）
\item $N$：齿轮比（通常$N > 1$，表示减速）
\end{itemize}

\textbf{高度齿轮传动}：$N \gg 1$（例如，$N = 100$或更大）

\subsection{高度齿轮传动的特性}

\textbf{从电机侧看}：
\begin{itemize}
\item 电机需要较小的力矩$\tau_m$就能产生较大的关节力矩$\tau_j = N\tau_m$
\item 电机可以高速旋转$\omega_m$，但关节速度$\omega_j = \omega_m/N$较慢
\item 电机惯性$I_m$被"放大"到关节侧：$I_j = N^2 I_m$
\end{itemize}

\textbf{从关节侧看}：
\begin{itemize}
\item 关节力矩$\tau_j$很大（相对于电机力矩）
\item 关节速度$\omega_j$较慢
\item 关节惯性$I_j$很大（$I_j = N^2 I_m$）
\item \textbf{很难被推动}（高阻抗）
\end{itemize}

\subsection{反向驱动性}

\textbf{反向驱动性（Backdrivability）}：
\begin{itemize}
\item 定义：从关节侧推动机器人，使其运动的难易程度
\item 高度齿轮传动的机器人：\textbf{很难反向驱动}
\item 原因：需要很大的力才能克服齿轮的摩擦和惯性
\end{itemize}

\textbf{数学分析}：

假设用户施加力$f_{\text{ext}}$在末端执行器上，产生的关节力矩为：
\begin{equation}
\tau_{\text{ext}} = J^T(\theta) f_{\text{ext}}
\end{equation}

要推动关节，需要克服：
\begin{enumerate}
\item 齿轮摩擦：$\tau_{\text{friction}}$
\item 关节惯性：$I_j \ddot{\theta}$
\item 电机惯性（放大后）：$N^2 I_m \ddot{\theta}$
\end{enumerate}

\textbf{总阻力}：
\begin{equation}
\tau_{\text{resistance}} = \tau_{\text{friction}} + (I_j + N^2 I_m)\ddot{\theta}
\end{equation}

由于$N^2 I_m$很大，$\tau_{\text{resistance}}$很大，因此很难推动。

\section{导纳控制的工作原理}

\subsection{导纳控制公式}

\textbf{导纳控制公式}：
\begin{equation}
M\ddot{x}_d + B\dot{x} + Kx = f_{\text{ext}} \tag{11.64-admittance}
\end{equation}

其中：
\begin{itemize}
\item $f_{\text{ext}}$：用户施加的力（由力传感器测量）
\item $x, \dot{x}$：当前位置和速度（由编码器和转速计测量）
\item $M, B, K$：虚拟质量、阻尼和刚度
\item $\ddot{x}_d$：期望加速度（计算得到）
\end{itemize}

\textbf{求解期望加速度}：
\begin{equation}
\ddot{x}_d = M^{-1}(f_{\text{ext}} - B\dot{x} - Kx) \tag{11.66}
\end{equation}

\subsection{导纳控制的实现步骤}

\textbf{步骤1：测量外力}
\begin{itemize}
\item 使用力-力矩传感器测量$f_{\text{ext}}$
\item 测量当前位置$x$和速度$\dot{x}$
\end{itemize}

\textbf{步骤2：计算期望加速度}
\begin{equation}
\ddot{x}_d = M^{-1}(f_{\text{ext}} - B\dot{x} - Kx)
\end{equation}

\textbf{步骤3：转换为关节加速度}
\begin{equation}
\ddot{\theta}_d = J^{\dagger}(\theta)(\ddot{x}_d - \dot{J}(\theta)\dot{\theta})
\end{equation}

\textbf{步骤4：计算关节力矩}
\begin{equation}
\tau = M(\theta)\ddot{\theta}_d + C(\theta, \dot{\theta})\dot{\theta} + g(\theta)
\end{equation}

\textbf{步骤5：执行控制}
\begin{itemize}
\item 电机产生力矩$\tau_m = \tau/N$（因为$\tau_j = N\tau_m$）
\item 关节产生加速度$\ddot{\theta}$
\item 末端执行器产生加速度$\ddot{x}$
\end{itemize}

\section{为什么高度齿轮传动适合模拟刚性环境？}

\subsection{刚性环境的特征}

\textbf{刚性环境}：
\begin{itemize}
\item 高刚度$K$（很大）
\item 小的位移$x$产生大的力$Kx$
\item 很难推动（高阻抗）
\item 位置变化很小，即使有外力
\end{itemize}

\textbf{数学表示}：
\begin{equation}
f = Kx \quad \Rightarrow \quad x = \frac{f}{K}
\end{equation}

如果$K$很大，即使$f$很大，$x$也很小。

\subsection{高度齿轮传动的优势}

\textbf{优势1：高阻抗特性}

\textbf{从关节侧看}：
\begin{itemize}
\item 高度齿轮传动的机器人本身就有高阻抗
\item 很难被推动（因为$N^2 I_m$很大）
\item 即使不控制，位置也相对稳定
\end{itemize}

\textbf{在导纳控制中}：
\begin{itemize}
\item 机器人主动响应力，产生运动
\item 但由于高阻抗，运动很小
\item 即使有外力$f_{\text{ext}}$，位置变化$x$也很小
\end{itemize}

\textbf{优势2：精确的位置控制}

\textbf{高齿轮比的优势}：
\begin{itemize}
\item 电机可以产生很大的关节力矩$\tau_j = N\tau_m$
\item 可以精确控制位置（即使有外力干扰）
\item 位置误差很小
\end{itemize}

\textbf{在导纳控制中}：
\begin{itemize}
\item 计算期望加速度$\ddot{x}_d$
\item 转换为关节力矩$\tau$
\item 高齿轮比使得可以产生很大的$\tau_j$，精确跟踪$\ddot{x}_d$
\item 结果：位置控制精确，即使有外力
\end{itemize}

\textbf{优势3：稳定性}

\textbf{高阻抗的稳定性}：
\begin{itemize}
\item 高阻抗意味着系统"硬"，不容易振荡
\item 即使有噪声或延迟，系统也相对稳定
\end{itemize}

\textbf{在导纳控制中}：
\begin{itemize}
\item 高齿轮比提供额外的"阻尼"（通过高惯性）
\item 减少振荡和不稳定
\item 即使虚拟刚度$K$很大，系统也稳定
\end{itemize}

\subsection{数学分析}

\textbf{导纳控制公式}：
\begin{equation}
M\ddot{x}_d + B\dot{x} + Kx = f_{\text{ext}}
\end{equation}

\textbf{对于刚性环境}，$K$很大，因此：
\begin{equation}
Kx \approx f_{\text{ext}} \quad \Rightarrow \quad x \approx \frac{f_{\text{ext}}}{K}
\end{equation}

\textbf{高度齿轮传动的效果}：

\textbf{1. 实际阻抗增加}：

由于高齿轮比，实际关节惯性$I_j = N^2 I_m$很大，这增加了系统的实际阻抗。

\textbf{有效刚度}：
\begin{equation}
K_{\text{effective}} = K + K_{\text{mechanical}}
\end{equation}

其中$K_{\text{mechanical}}$是由高齿轮比产生的机械刚度。

\textbf{2. 位置误差减小}：

由于$K_{\text{effective}}$很大，位置误差$x$很小：
\begin{equation}
x = \frac{f_{\text{ext}}}{K_{\text{effective}}} \approx \frac{f_{\text{ext}}}{K + K_{\text{mechanical}}}
\end{equation}

\textbf{3. 控制精度提高}：

高齿轮比使得可以产生很大的关节力矩，精确跟踪期望加速度：
\begin{equation}
\tau_j = N\tau_m = N \cdot \frac{\tau_{\text{computed}}}{N} = \tau_{\text{computed}}
\end{equation}

（实际上，由于齿轮效率$\eta < 1$，$\tau_j = \eta N\tau_m$，但仍然很大）

\section{与低质量环境模拟的对比}

\subsection{低质量环境的特征}

\textbf{低质量环境}：
\begin{itemize}
\item 小质量$M$（很小）
\item 小的力$f_{\text{ext}}$产生大的加速度$\ddot{x}_d = M^{-1}f_{\text{ext}}$
\item 很容易推动（低阻抗）
\item 位置变化很大，即使有小的外力
\end{itemize}

\textbf{数学表示}：
\begin{equation}
M\ddot{x}_d = f_{\text{ext}} \quad \Rightarrow \quad \ddot{x}_d = \frac{f_{\text{ext}}}{M}
\end{equation}

如果$M$很小，即使$f_{\text{ext}}$很小，$\ddot{x}_d$也很大。

\subsection{为什么导纳控制难以模拟低质量环境？}

\textbf{问题1：需要大的加速度}

\textbf{对于低质量环境}：
\begin{equation}
\ddot{x}_d = M^{-1}(f_{\text{ext}} - B\dot{x} - Kx) \approx M^{-1}f_{\text{ext}}
\end{equation}

如果$M$很小，$\ddot{x}_d$很大。

\textbf{对于高度齿轮传动的机器人}：
\begin{itemize}
\item 关节速度$\omega_j = \omega_m/N$较慢
\item 难以产生大的加速度$\ddot{x}_d$
\item 因为$\ddot{x}_d = J(\theta)\ddot{\theta} + \dot{J}(\theta)\dot{\theta}$，而$\ddot{\theta}$受限于$\omega_m$和$N$
\end{itemize}

\textbf{问题2：有效增益很大}

\textbf{有效增益}：
\begin{equation}
K_{\text{effective}} = \frac{f_{\text{ext}}}{\Delta x} = \frac{M\ddot{x}_d + B\dot{x} + Kx}{\Delta x}
\end{equation}

如果$M$很小，$\ddot{x}_d$很大，$K_{\text{effective}}$很大。

\textbf{在高增益下}：
\begin{itemize}
\item 系统容易不稳定
\item 对噪声和延迟敏感
\item 可能振荡
\end{itemize}

\textbf{问题3：反向驱动性要求}

\textbf{低质量环境需要}：
\begin{itemize}
\item 机器人容易被推动（低阻抗）
\item 小的力产生大的运动
\end{itemize}

\textbf{高度齿轮传动的机器人}：
\begin{itemize}
\item 很难被推动（高阻抗）
\item 不适合模拟低质量环境
\end{itemize}

\subsection{为什么高度齿轮传动不适合模拟低质量环境？}

\textbf{原因1：高阻抗}

\textbf{高度齿轮传动的机器人}：
\begin{itemize}
\item 从关节侧看，阻抗很高（$N^2 I_m$很大）
\item 很难被推动
\item 不适合模拟低阻抗（低质量）环境
\end{itemize}

\textbf{原因2：速度限制}

\textbf{高度齿轮传动的机器人}：
\begin{itemize}
\item 关节速度$\omega_j = \omega_m/N$较慢
\item 难以产生大的加速度$\ddot{x}_d$
\item 不适合模拟低质量环境（需要大的加速度）
\end{itemize}

\textbf{原因3：控制困难}

\textbf{在导纳控制中}：
\begin{itemize}
\item 需要快速响应外力
\item 需要产生大的加速度
\item 高度齿轮传动的机器人难以做到这一点
\end{itemize}

\section{为什么高度齿轮传动适合模拟刚性环境？}

\subsection{总结}

\textbf{高度齿轮传动的机器人的特点}：
\begin{enumerate}
\item \textbf{高阻抗}：很难被推动（$N^2 I_m$很大）
\item \textbf{精确控制}：可以产生很大的关节力矩，精确控制位置
\item \textbf{稳定性}：高阻抗提供额外的阻尼，减少振荡
\end{enumerate}

\textbf{刚性环境的特征}：
\begin{enumerate}
\item \textbf{高刚度}：小的位移产生大的力
\item \textbf{高阻抗}：很难推动
\item \textbf{位置稳定}：即使有外力，位置变化也很小
\end{enumerate}

\textbf{匹配性}：
\begin{itemize}
\item 高度齿轮传动的机器人的高阻抗特性与刚性环境的高阻抗特性匹配
\item 可以精确控制位置，即使有外力干扰
\item 系统稳定，不容易振荡
\end{itemize}

\subsection{在导纳控制中的表现}

\textbf{导纳控制公式}：
\begin{equation}
M\ddot{x}_d + B\dot{x} + Kx = f_{\text{ext}}
\end{equation}

\textbf{对于刚性环境}，$K$很大，因此：
\begin{equation}
x \approx \frac{f_{\text{ext}}}{K}
\end{equation}

\textbf{高度齿轮传动的效果}：
\begin{enumerate}
\item \textbf{实际阻抗增加}：$K_{\text{effective}} = K + K_{\text{mechanical}}$很大
\item \textbf{位置误差减小}：$x = f_{\text{ext}}/K_{\text{effective}}$很小
\item \textbf{控制精度提高}：可以产生很大的关节力矩，精确跟踪期望加速度
\item \textbf{稳定性提高}：高阻抗提供额外的阻尼，减少振荡
\end{enumerate}

\section{实际例子}

\subsection{例子1：刚性环境模拟}

\textbf{场景}：
\begin{itemize}
\item 虚拟刚度：$K = 10000$ N/m（很大）
\item 用户施加力：$f_{\text{ext}} = 10$ N
\item 期望位置变化：$x = f_{\text{ext}}/K = 0.001$ m（1毫米）
\end{itemize}

\textbf{高度齿轮传动的机器人}：
\begin{itemize}
\item 高齿轮比：$N = 100$
\item 关节惯性：$I_j = N^2 I_m = 10000 I_m$（很大）
\item 可以产生很大的关节力矩：$\tau_j = N\tau_m = 100\tau_m$
\item 精确控制位置，位置误差很小（约1毫米）
\end{itemize}

\textbf{结果}：成功模拟刚性环境。

\subsection{例子2：低质量环境模拟（困难）}

\textbf{场景}：
\begin{itemize}
\item 虚拟质量：$M = 0.1$ kg（很小）
\item 用户施加力：$f_{\text{ext}} = 1$ N
\item 期望加速度：$\ddot{x}_d = f_{\text{ext}}/M = 10$ m/s²（很大）
\end{itemize}

\textbf{高度齿轮传动的机器人}：
\begin{itemize}
\item 关节速度较慢：$\omega_j = \omega_m/N = \omega_m/100$
\item 难以产生大的加速度：$\ddot{x}_d = J(\theta)\ddot{\theta} + \dot{J}(\theta)\dot{\theta}$
\item 由于$\ddot{\theta}$受限于$\omega_m$和$N$，$\ddot{x}_d$难以达到10 m/s²
\end{itemize}

\textbf{结果}：难以模拟低质量环境。

\section{总结}

\subsection{核心结论}

\begin{enumerate}
\item \textbf{高度齿轮传动的特点}：
\begin{itemize}
\item 高阻抗（很难被推动）
\item 精确的位置控制能力
\item 稳定性好
\end{itemize}

\item \textbf{刚性环境的特征}：
\begin{itemize}
\item 高刚度（小的位移产生大的力）
\item 高阻抗（很难推动）
\item 位置稳定
\end{itemize}

\item \textbf{匹配性}：
\begin{itemize}
\item 高度齿轮传动的机器人的高阻抗特性与刚性环境的高阻抗特性匹配
\item 在导纳控制中，可以很好地模拟刚性环境
\item 但难以模拟低质量环境（需要低阻抗和大的加速度）
\end{itemize}

\item \textbf{导纳控制的优势}：
\begin{itemize}
\item 感测力，产生运动
\item 高度齿轮传动提供额外的机械阻抗
\item 使得虚拟刚度$K$可以设置得很大，而不导致不稳定
\item 位置控制精确，即使有外力干扰
\end{itemize}
\end{enumerate}

\subsection{关键公式}

\textbf{导纳控制}：
\begin{equation}
M\ddot{x}_d + B\dot{x} + Kx = f_{\text{ext}}
\end{equation}

\textbf{高度齿轮传动的效果}：
\begin{equation}
K_{\text{effective}} = K + K_{\text{mechanical}} \quad \Rightarrow \quad x = \frac{f_{\text{ext}}}{K_{\text{effective}}}
\end{equation}

\textbf{结果}：位置误差$x$很小，成功模拟刚性环境。

\end{document}

